% Generated draft: requires mathematical and source review.
\documentclass[11pt]{article}
\usepackage[margin=25mm]{geometry}
\usepackage{fontspec}
\setmainfont{Latin Modern Roman}
\usepackage{amsmath,amssymb,mathtools,unicode-math}
\setmathfont{Latin Modern Math}
\usepackage{enumitem,needspace,xurl,hyperref}
\setlength{\parindent}{0pt}
\setlength{\parskip}{5pt}
\setlength{\emergencystretch}{3em}
\newcommand{\N}{\mathbb{N}}
\newcommand{\Z}{\mathbb{Z}}
\newcommand{\Q}{\mathbb{Q}}
\newcommand{\R}{\mathbb{R}}
\newcommand{\C}{\mathbb{C}}
\newcommand{\F}{\mathbb{F}}
\newcommand{\A}{\mathbb{A}}
\newcommand{\PP}{\mathbb{P}}
\newcommand{\E}{\mathbb{E}}
\newcommand{\eps}{\varepsilon}
\newcommand{\dd}{\,\mathrm{d}}
\begin{document}
\section*{UnsolvedMath ID: 701248}
\textbf{Source problem code:} catalogue entry
\subsection{Definitions and mathematical statement}
For a group \(G\) with identity \(1_G\), define the order of \(g\in G\) by
\[
\operatorname{ord}_G(g)=\min\{n\in\Z_{\geq 1}:g^n=1_G\},
\]
with \(\operatorname{ord}_G(g)=\infty\) if this set is empty. The spectrum of \(G\) is
\[
\omega(G)=\{\operatorname{ord}_G(g):g\in G\}.
\]
A group is \emph{locally finite} if every subgroup generated by finitely many elements is finite.

Is it true that, for every group \(G\),
\[
\omega(G)=\{1,2,3,4,5,6\}
\quad\Longrightarrow\quad
|\langle S\rangle|<\infty\text{ for every finite subset }S\subseteq G?
\]
The spectrum hypothesis requires every element to have one of the six listed orders and requires each listed order to occur. No finiteness or finite generation assumption on \(G\) is imposed.
\subsection{Short English statement}
Must every group whose element orders are exactly \(1,2,3,4,5,6\) have all its finitely generated subgroups finite? This question is attributed to V. D. Mazurov (2006). The supplied registry labels it open; that status has not been independently verified.
\subsection{Sources}
\begin{itemize}
\item Local UnsolvedMath registry, permanent ID \(701248\), exact source problem code \texttt{catalogue entry}, title \emph{Kourovka Notebook Problem 16.56}. The supplied record labels the problem open.
\item V. D. Mazurov, Problem 16.56, \emph{The Kourovka Notebook (Unsolved Problems in Group Theory)}, 16th issue (2006), as attributed in the supplied record.
\item The supplied background also identifies the 21st issue edition (2026), arXiv identifier \texttt{1401.0300v48}. The supplied collection description names E. I. Khukhro and V. D. Mazurov as editors. These bibliographic details are reproduced from the record and have not been independently verified.
\end{itemize}
\end{document}
